% ============================================================
% reviewer-annotation-spec.tex
% Cosmetic, review-only markup: highlights reviewer comments and their
% resolutions directly on the manuscript as Word-style margin notes.
% Reviewer 1 = orange, Reviewer 2 = blue.
%
% This file is \input only by paper-annotated.tex. The clean manuscript
% (paper.tex + sections/*.tex) does not load it and is unchanged.
%
% To produce a clean (markup-free) build, load changes with [final]:
%   \usepackage[final]{changes}
% ============================================================

% Float-free margin notes (Word-style side comments). marginnote is used
% instead of todonotes so that comments next to MDPI's [H] tables/figures
% are never dropped ("Float(s) lost").
\usepackage{marginnote}
\renewcommand*{\marginfont}{\scriptsize\raggedright}

% Per-author change tracking with margin comments
\usepackage[
  markup=underlined,
  authormarkup=superscript,
  commentmarkup=default
]{changes}

% Reviewers as "authors" with distinct colours
\definechangesauthor[name={Reviewer 1}, color=orange]{Rone}
\definechangesauthor[name={Reviewer 2}, color=blue]{Rtwo}

% Math/macro-safe markup: highlighted spans only recolour the text (the
% per-author colour is applied automatically), so they may safely contain
% inline math ($...$), \emph, \texttt, \ref, etc. without soul/ulem breakage.
% "authorcolor" is set by changes to the current reviewer's colour before the
% markup expands; wrapping in \textcolor keeps per-reviewer colour while being
% safe for inline math and macros inside the span.
\sethighlightmarkup{\textcolor{authorcolor}{#1}}
\setaddedmarkup{\textcolor{authorcolor}{#1}}

% Route reviewer comments to float-free margin notes (coloured per author).
\setcommentmarkup{\marginnote{\textcolor{authorcolor}{[#1]}}}

% ===== Convenience macros (see docs/overleaf-comments.md) =====
% <comment> is a short "Cxx: comment -> resolution" string shown in the margin.

% --- Reviewer 1 (orange) ---
\newcommand{\RoneAdd}[2][]{\added[id=Rone, comment={#1}]{#2}}
\newcommand{\RoneDel}[2][]{\deleted[id=Rone, comment={#1}]{#2}}
\newcommand{\RoneRep}[3][]{\replaced[id=Rone, comment={#1}]{#3}{#2}}
\newcommand{\RoneHl}[2][]{\highlight[id=Rone, comment={#1}]{#2}}
\newcommand{\RoneCom}[1]{\comment[id=Rone]{#1}}

% --- Reviewer 2 (blue) ---
\newcommand{\RtwoAdd}[2][]{\added[id=Rtwo, comment={#1}]{#2}}
\newcommand{\RtwoDel}[2][]{\deleted[id=Rtwo, comment={#1}]{#2}}
\newcommand{\RtwoRep}[3][]{\replaced[id=Rtwo, comment={#1}]{#3}{#2}}
\newcommand{\RtwoHl}[2][]{\highlight[id=Rtwo, comment={#1}]{#2}}
\newcommand{\RtwoCom}[1]{\comment[id=Rtwo]{#1}}
